% !TEX root = ../main.tex
% ============================================================================
% Conclusion & Future Work
% ============================================================================
\chapter{Conclusion and Future Work}
\label{ch:conclusion}

\section{Summary}
\label{sec:concl_summary}
This thesis asked whether knowledge distillation can make a small HiVT trajectory
predictor competitive with a much larger one without paying a calibration cost. The
answer is yes, but only with the right objective. The contribution is a
\emph{permutation-invariant, distribution-matching} distillation loss for
Laplace-mixture predictors, together with the analysis and experiments that explain
why the obvious alternatives fail.

% TODO: tighten into 1-2 paragraphs once the surrounding chapters are final.

\section{Contributions}
\label{sec:concl_contrib}
The main contributions, by chapter, are:

\begin{itemize}[leftmargin=1.4em]
  \item \textbf{A diagnosis of why index-aligned distillation is ill-posed for
        mixture predictors} (Chapter~\ref{ch:method}). The winner-takes-all
        training of HiVT leaves its mode \emph{slots} unconstrained, so a teacher
        and an independently trained student number their modes under a consistent
        but non-identity permutation. Pairing modes by index is optimal in $0\%$ of
        scenes and mismatches modes by $\approx 4\times$ the optimal-assignment
        distance (Figure~\ref{fig:perm}).
  \item \textbf{A permutation-invariant distillation objective} that scores the
        teacher's modes under the student's full mixture density, requiring no mode
        correspondence and remaining well-defined when teacher and student differ in
        mode count (Section~\ref{ssec:method_pi}).
  \item \textbf{An analysis of the calibration failure of mean-target
        distillation} (v1) and its fix (v2). Discarding the teacher's scales rewards
        unbounded scale shrinkage; restoring them through a distribution-matching
        (Monte-Carlo) target removes the pathology, and v1 is recovered exactly as
        the zero-variance limit of v2 (Sections~\ref{ssec:method_calib},
        \ref{ssec:method_v2}).
  \item \textbf{Empirical confirmation at two student widths} (Chapter~\ref{ch:results}).
        On HiVT-32, v2 matches v1's geometric gain (minFDE $-9.2\%$) while repairing
        calibration to \emph{better than the teacher}; on HiVT-16 the gain is larger
        still (minFDE $-22.7\%$) with calibration also improving. Distillation buys
        roughly a full size-class of accuracy at zero inference cost, and the smaller
        the student the larger the relative benefit.
  \item \textbf{An efficiency characterisation} clarifying that the unconditional
        dividend of a smaller model is \emph{memory} (up to $6.3\times$ less GPU
        memory), while the latency win is real only on CPU under modest batching
        (Section~\ref{sec:res_eff}).
\end{itemize}

\section{Limitations}
\label{sec:concl_limits}
Several limitations bound the claims of this work and motivate the directions below.

\begin{itemize}[leftmargin=1.4em]
  \item \textbf{Single dataset and backbone.} All results are on Argoverse~1 with
        the HiVT architecture; generality to other datasets (Argoverse~2,
        nuScenes, Waymo) and other mixture-based predictors is untested.
  \item \textbf{Limited seed counts.} The HiVT-16 A/B is confirmed over multiple
        seeds (no~KD $n{=}2$, v2 $n{=}3$; Section~\ref{ssec:res_seeds}), where the
        gain is reproducible and well outside the seed spread, but the HiVT-32 A/Bs
        remain single training runs. Most effects are far larger than plausible
        run-to-run noise; the smaller calibration differences at HiVT-32 would be
        firmer with seed replication.
  \item \textbf{Proxy-based hyperparameter selection.} The $\lambda$ and recipe
        choices are ranked on a $25\%$-data, short-schedule proxy. The proxy ranked
        correctly and even \emph{under}-stated the full-data gain here, but it is
        not guaranteed to transfer in general.
  \item \textbf{One teacher.} Only HiVT-128 is used as the teacher; the effect of
        teacher capacity, or of multi-teacher / self-distillation, is not studied.
\end{itemize}

\section{Future Work}
\label{sec:concl_future}

\begin{itemize}[leftmargin=1.4em]
  \item \textbf{Multi-seed variance and training stability.} Extend the HiVT-16
        seed study (Section~\ref{ssec:res_seeds}) to the HiVT-32 A/Bs, reporting
        mean$\pm$spread to harden the calibration claims
        (Section~\ref{ssec:res_v2}). A related thread is the observation that the
        from-scratch HiVT-16 occasionally collapses into a degenerate optimum
        while none of the distilled runs did: a larger seed sweep could test
        whether distillation reliably \emph{stabilises} training away from that
        basin.
  \item \textbf{Other backbones and datasets.} Apply the permutation-invariant v2
        loss to other mixture predictors and to Argoverse~2 / nuScenes, testing
        whether the ``calibration-preserving, larger-gain-for-smaller-students''
        pattern is architecture- and dataset-independent.
  \item \textbf{Teacher capacity and ensembles.} Vary the teacher width and try
        ensemble or multi-teacher distillation, where the soft-target set of
        Section~\ref{ssec:method_pi} extends naturally to a union of mixtures.
  \item \textbf{A scale-aware $\lambda$ schedule.} Because $\lambda$ is not
        comparable across v1 and v2, and the v2 optimum drifts with student width
        ($0.5$ at HiVT-32, $1.0$ at HiVT-16), an automatic balancing of the task and
        distillation terms (e.g.\ gradient-norm or uncertainty weighting) would
        remove a manual sweep.
  \item \textbf{Closed-loop / downstream evaluation.} The motivation for calibrated
        uncertainty is a downstream planner; evaluating the distilled student inside
        a planning or closed-loop simulation would test whether the calibration
        gains translate into better decisions, not just better likelihoods.
  \item \textbf{On-device deployment.} Given that the dividend is memory and that
        CPU is fastest at batch size~$1$, an INT8/quantised, CPU-targeted deployment
        of the distilled student is a natural next step
        (Section~\ref{sec:res_eff}).
\end{itemize}
